\subsection{A negative word trace on an actual rank-four simple module}
The rank-three coefficient kernels do not extend, with the same
strict positivity requirement, to every simple module of the full
rank-four nonstandard Hecke algebra. The obstruction is a trace of
a named word and is therefore independent of the choice of basis.
It concerns an additional simple module outside the two-row quotient.

\begin{finding}[\fid{F57}: an actual rank-four word-trace obstruction]
\label{f:rankfournegativetrace}
Put \(\delta=v+v^{-1}\), \(D=\delta^2\), and
\[
 N=S'\check M_{(3,1)},\qquad
 N^\#=\text{the }\Theta^{\mathrm{op}}\text{-contragredient of }N.
\]
The latter is the additional five-dimensional simple module listed
in IV, \S11.8, p.~67. For the named word
\[
 \omega=P_2P_3P_2P_3P_2P_3P_1P_2P_1P_2P_1
\]
one has
\begin{equation}\label{eq:rankfournegativepolynomial}
\begin{split}
 \operatorname{Tr}_{N^\#}(\omega)=p(D)
 ={}&-D^9+20D^8-163D^7+734D^6-2048D^5\\
    &+3712D^4-4352D^3+3072D^2-1024D.
\end{split}
\end{equation}
Its Laurent coefficients at \(v^{18}\) and \(v^{-18}\) are \(-1\).
The module \(N^\#\) is an actual direct summand of the restriction
of \(M_{(2,1,1)}\otimes M_{(3,1)}\), hence occurs in the actual
nonstandard tensor representation when
\(\dim V\ge3,\ \dim W\ge2\), or with the factors interchanged.

Consequently no basis of this specified generic simple module
makes all three named \(P_i\) matrices coefficientwise nonnegative
Laurent polynomials. In particular there is no effective finite
graded-matrix-kernel realization with this Grothendieck module
and these named generators.

For any fixed \(q>1\), neither
\[
 \pm p\big((q^{n/2}+q^{-n/2})^2\big)
\]
is the power-trace sequence of an invertible endomorphism of a
finite-dimensional vector space for all sufficiently large \(n\).
A common Tate twist does not repair this. A bounded complex
with both cohomological parities does realize the scalar sequence
when \(q\) is a prime power.

The negative \(P\)-word occurs on the extra full-algebra
\(N^\#\) block. It does not exclude a \(P_i\)-only effective
kernel model for the two-row quotient, nor does it contradict
the signed/factored positivity stated in IV.
\end{finding}

\begin{proof}
\par\smallskip\noindent\textbf{The exact source input.}\quad
IV already constructs the two-row tensor basis through rank four:
see \S1.5, p.~13, Conjecture~19.1, pp.~130--131, and
Example~19.3, pp.~133--135.
Its stated coefficient property is signed and factored:
Remark~11.19, p.~67, and Example~19.3, p.~134, allow signs and
\(c_0=[2]^2-4\).
This is distinct from strict nonnegative Laurent matrices.
The canonical cellular basis of the algebra itself remains
unconstructed there at rank four; see \S19.1, p.~136
\cite{IV}.

Use the explicit \(Q_i\) matrices in Figure~22, p.~135,
with \(Q_i=D-P_i\).
This normalization follows from \S11.1, p.~59.
It also avoids propagating the exponent inconsistency between
that section's preliminary \(C'_s,C_s\) expressions and
\S3.1: the displayed matrices themselves satisfy \(Q_i^2=DQ_i\).

Set \(a=D-2,\ c=D-4\). Order the Figure~22 states as its
two upper states from left to right, followed by its three
lower states from left to right. Their descent sets are
\[
 \{2\},\quad\{1,3\},\quad\{1\},\quad\{2\},\quad\{3\}.
\]
Use right row matrices: entry \((r,s)\) is the coefficient of
state \(s\) in the action on state \(r\). The source's descent
and arrow conventions then give
\begin{equation}\label{eq:rankfourfivecellmatrices}
\begin{gathered}
Q_1=\begin{pmatrix}
0&-\delta&c&0&0\\
0&D&0&0&0\\
0&0&D&0&0\\
0&0&a&0&0\\
0&-\delta&0&0&0
\end{pmatrix},\qquad
Q_2=\begin{pmatrix}
D&0&0&0&0\\
-\delta&0&0&0&0\\
0&0&0&a&0\\
0&0&0&D&0\\
a&0&0&-a&0
\end{pmatrix},\\[2mm]
Q_3=\begin{pmatrix}
0&-\delta&0&0&c\\
0&D&0&0&0\\
0&-\delta&0&0&0\\
0&-\delta&0&0&-2\\
0&0&0&0&D
\end{pmatrix}.
\end{gathered}
\end{equation}
They are matrices of the actual \(N=S'\check M_{(3,1)}\)
in Example~19.3, not an abstract representation identified
only by a dimension count.

\par\smallskip\noindent\textbf{An explicit product certificate.}\quad
Write
\[
 w=(1,2,1,2,1,3,2,3,2,3,2),\qquad
 Q_w=Q_1Q_2Q_1Q_2Q_1Q_3Q_2Q_3Q_2Q_3Q_2.
\]
Let
\[
 u=(-2\delta a,D^2,-\delta D,-\delta a,-\delta D)^t,
 \qquad e_2^t=(0,1,0,0,0).
\]
Equation~\eqref{eq:rankfourfivecellmatrices} gives
\(Q_1Q_3=ue_2^t\). After rotating inside the trace, put
\[
 A=Q_2Q_3Q_2Q_3Q_2Q_1Q_2Q_1Q_2.
\]
Successive multiplication of the row \(e_2^t\) by the
displayed matrices gives
\[
 e_2^tA=(r_1,0,0,r_4,0),
\]
where
\[
 r_1=-\delta^5(D^4-9D^3+31D^2-48D+32)
\]
and
\[
 r_4=\delta(D^7-16D^6+109D^5-418D^4
                 +992D^3-1472D^2+1280D-512).
\]
It follows that
\begin{align*}
 \operatorname{Tr}(Q_w)
 &=e_2^tAu=-\delta a(2r_1+r_4)\\
 &=-D(D-2)(D^7-18D^6+127D^5-480D^4\\
 &\hspace{37mm}+1088D^3-1536D^2+1280D-512),
\end{align*}
which expands to \(p(D)\).
For completeness, the only nonzero diagonal entries of
the full product are
\[
\begin{aligned}
(Q_w)_{11}={}&D^{10}-20D^9+180D^8-951D^7+3246D^6\\
&-7440D^5+11520D^4-11712D^3+7168D^2-2048D,\\
(Q_w)_{44}={}&-D^{10}+19D^9-160D^8+788D^7-2512D^6\\
&+5392D^5-7808D^4+7360D^3-4096D^2+1024D.
\end{aligned}
\]
Their sum again gives
\eqref{eq:rankfournegativepolynomial}.

\par\smallskip\noindent\textbf{The source dual and actual occurrence.}\quad
IV, Proposition~11.6, p.~60, supplies \(\Theta(P_i)=Q_i\).
Section~11.4, pp.~62--63, defines \(\#\) using the
anti-involution \(\Theta^{\mathrm{op}}\).
Thus the dual row matrices of \(P_i\) on \(N^\#\) are \(Q_i^t\).
The reversal in
\[
 \omega=P_{w^{\mathrm{rev}}}
\]
is essential to the direct matrix identity
\[
 \rho_{N^\#}(\omega)=(Q_w)^t.
\]
Taking traces proves
\eqref{eq:rankfournegativepolynomial} with the actual
source convention.

The generic direct-summand inclusion
\[
 N\subset S^2M_{(3,1)}
      \subset M_{(3,1)}\otimes M_{(3,1)}
\]
uses IV, (99), p.~63, and the symmetric/exterior splitting.
Apply the exact contragredient functor.
Propositions~11.11--11.12, p.~63, especially (101) and (105),
give
\[
 N^\#\subset
 \operatorname{Res}_{\check H_4}
   \big(M_{(2,1,1)}\otimes M_{(3,1)}\big)
\]
as a direct summand.
Ordinary Schur--Weyl duality puts this restriction in
\((V\otimes W)^{\otimes4}\) for \(\dim V\ge3,\ \dim W\ge2\).
The identification with the actual nonstandard action is
IV, Theorem~12.1, (115), p.~68, and (116)--(117), p.~69.
For dimensions \(3,2\), the ordinary left factors have
dimensions \(3,3\), so this particular isotypic component
contains nine copies of \(N^\#\).
Remark~12.2, p.~68, qualifies the source's \(X/X^*\)
convention. Consistent dualization transposes matrices and
reverses words; it preserves this trace obstruction.

The full-algebra list in IV, \S11.8, p.~67, includes
the additional five-dimensional \(N^\#\).
Theorem~13.4, pp.~74--75, gives the two-row list without it.
The latter squared dimensions sum to
\[
 1+4+9+36+25+4+9+1=89;
\]
the full-algebra dimension is \(89+25=114\).

\par\smallskip\noindent\textbf{Effective-kernel obstruction.}\quad
Expanding in \(v\) gives
\begin{equation}\label{eq:rankfournegativeLaurent}
\begin{split}
p(D)={}&-(v^{18}+v^{-18})+2(v^{16}+v^{-16})
 +4(v^{14}+v^{-14})\\
&+36(v^{12}+v^{-12})+67(v^{10}+v^{-10})
 +176(v^8+v^{-8})\\
&+297(v^6+v^{-6})+476(v^4+v^{-4})
 +657(v^2+v^{-2})+668.
\end{split}
\end{equation}
In particular its two outer coefficients are negative.
If every \(P_i\) had entries in
\(\mathbb N[v,v^{-1}]\), matrix multiplication and trace
would force the same property for every word trace.
Similarity cannot change this trace, so the obstruction
holds for every basis of the specified module.
Effective finite graded-vector-space kernels have exactly
this nonnegative matrix-character property.

There is also a pointwise obstruction:
\[
 p((3+3^{-1})^2)
 =-\frac{99834597608550400}{387420489}<0.
\]
On this displayed specialization, no real basis makes all
three \(P_i\) matrices entrywise nonnegative.
No irreducibility assertion at \(v=3\) is needed.

On a larger module the same conclusion follows if one also
requires an effective kernel projector onto this isotypic
component: its inserted word trace is a positive multiplicity
times \(p(D)\).
Effectivity of that projector is an additional premise;
an arbitrary positive basis of a larger tensor module
need not make its direct-summand projectors effective.

\par\smallskip\noindent\textbf{Finite power traces, parity, and a derived witness.}\quad
At \(v=q^{n/2}\),
\eqref{eq:rankfournegativeLaurent} is a finite exponential
sum with distinct roots \(q^{e/2}\) and the displayed
integer coefficients.
A genuine finite power trace
\(\operatorname{Tr}(F^n;H)\), with \(F\) invertible,
has nonnegative integer multiplicities after equal
eigenvalues are grouped. Jordan blocks do not change this
trace.
Subtracting the two expressions and applying the finite
Vandermonde independence in \fid{F02} contradicts the
coefficient \(-1\) at \(q^9\).
The same proof holds if equality is assumed only for
all sufficiently large \(n\).
The inverse parameter convention gives the same polynomial.

An overall parity shift cannot help:
\(p(D)\) has negative outer coefficients while
\(-p(D)\) has, for example, constant coefficient \(-668\).
A fixed Tate twist multiplies all distinct roots by one
nonzero scalar and leaves these signs unchanged.

There is nevertheless an explicit bounded Tate complex
realizing the scalar trace. For prime-power \(q\), choose
\(\ell\nmid q\) and let
\[
 (b_1,\ldots,b_8)=(657,476,297,176,67,36,4,2),
\]
\[
 E_+=\mathbb Q_\ell^{\,668}
   \oplus\bigoplus_{k=1}^{8}
      \big(\mathbb Q_\ell(-k)\oplus
           \mathbb Q_\ell(k)\big)^{\oplus b_k},
 \qquad
 E_-=\mathbb Q_\ell(-9)\oplus\mathbb Q_\ell(9).
\]
Then \(\dim E_+=4098\), and
\[
 E_+\oplus E_-[-1]
\]
has Frobenius Euler trace \(p(D)\) at every extension degree.
This proves the exact two-parity escape for this scalar
sequence; it supplies no compatible relation maps for
the original word action.
\end{proof}

The calculation is an obstruction to the strict effective
kernel extension on a specified full-algebra simple module.
It does not refute IV's signed/factored positivity, or
exclude alternating derived traces and inserted-operator
traces. On the original two-row \(N\), it is a \(Q\)-word
trace, and \(Q_i=D-P_i\) uses subtraction. Thus the
\(P_i\)-only two-row route remains outside this result.
The source matrices, product diagonals, rank-one reduction,
dual reversal, Laurent expansion and exact specialization
are also checked in the saved exact certificate; the proof
does not infer an all-word claim from a finite search.

% Separate paragraph for Part III.
\paragraph{\fid{F57}'s trace obstruction in Grothendieck's language.}
Finite graded coefficient spaces are representations of
a split torus. Their effective character cone is
\(\mathbb N[v,v^{-1}]\); duality reverses the character,
and tensor product multiplies characters.
The representation framework is SGA~3,
I, 4.7.2--3, with cote~90, scan~17 \cite{SGA3,Archive}.
This is action language, not a new scheme-group realization
of the noncommutative algebra.
SGA~6, V, \S\S2--3 and cote~133, scan~18 supply the
representation and \(\lambda/\sigma\) vocabulary
\cite{SGA6,Archive}; no \(\lambda\)-ring structure on the
noncommutative word algebra is asserted.
The actual module, its dual, and its endomorphism matrices
are affine module data in EGA~I, Theorem~1.3.7, p.~86,
Corollary~1.3.8, p.~87, and Theorem~1.4.1,
pp.~90--92 \cite{EGA}.
The new calculation shows that the specified word trace
lies outside the effective cone with either overall sign.
It remains a valid virtual character, as the explicit
bounded Tate complex demonstrates. Its underlying bounded
coefficient complex fits SGA~6, I, Definition~4.7,
pp.~121--122, and Example~4.8, p.~122 \cite{SGA6};
the Frobenius action is additional specified coefficient data.
The grading-to-Frobenius comparison and an effective
summand projector remain extra hypotheses.
This is an effectivity
and comparison obstruction, not a failure of a
Riemann-hypothesis theorem or a source claim that every
nonstandard action matrix was strictly positive.
